%% 毕业设计小结

\begin{designsummary}
本次毕业设计围绕大语言模型对齐中奖励函数动态更新的稳定性问题展开，主要研究如何在逆强化学习框架下约束奖励函数的变化幅度，并把理论分析转化为可实现的训练机制。从最初阅读RLHF和IRL方面的文献、到方法推导、实验实施、结果分析和论文撰写，整体工作让我对人工智能对齐这一类问题形成了相对系统的认识。

理论部分的工作要求我把强化学习、PPO、GRPO、RLHF、DPO和逆强化学习这些方法重新摆到同一个问题背景下来比较。过去对它们的理解多停留在概念层面，本课题则需要分析固定奖励模型下的分布偏移、奖励过优化以及IRL中奖励—策略交替优化的稳定性。讨论这些问题时我逐渐倾向于一个判断：大语言模型对齐的关键并不在于挑选某种策略优化算法，而在于策略不断变化时奖励信号是否仍然可靠、可控。

从理论推导走到实际训练系统的距离比预想的大。第三章给出的改进下界从数学上指明了相邻奖励函数变化幅度需要被控制，但其中的最大逐点变化和常数项在真实语言模型训练里并不容易直接计算。为了让结论能落到训练流程，本文用批次RMSD作为奖励变化的代理量，并配合自适应约束系数的调节机制。这一步让我注意到形式上正确的理论结论之外，还要考虑结论是否能被可计算、可调节的工程接口承接。

实验部分基于SLIME后训练框架，完成了TL;DR摘要生成与UltraFeedback通用指令跟随两组实验，期间接触到大语言模型后训练里的数据准备、策略采样、奖励更新、检查点管理和自动评测等环节。实际操作中遇到的问题不算少，包括训练参数的取值、不同随机种子下的结果波动以及多轮迭代结果的整理对齐。反复调试之后，我对大模型训练实验的复杂度有了更直接的判断：实验结论的说服力，建立在清晰的对比设置和能够解释训练动态的中间指标之上。论文撰写阶段的体会与此类似。一个方法能否讲清楚，并不只取决于公式是否正确，还取决于动机、推导、实验设计和结果分析能否串成连贯的论证。整理章节结构、统一符号、修订公式和图表的工作虽然繁琐，但直接决定论文能不能让人读下去。

本次工作仍存在明显局限。受限于时间和计算资源，实验主要覆盖两个任务和有限的模型规模，没有在更大模型、复杂多轮对话或真实在线反馈场景下验证；本文使用的也仅是标量奖励和批次级RMSD约束，对多维偏好、细粒度奖励变化以及更紧的理论误差界的处理尚不充分。这些是后续可以继续推进的方向。
\end{designsummary}
